% Agent 23: Pages 430--431 (PDF pp.44--45)
% Section 5.6 (cont.) -- energy conservation in radial structure, F formula,
%   manipulation of conservation laws, final radial structure equations
% Section 5.7 -- Equations of Vertical Structure (beginning)
% Equations (5.6.12)--(5.6.14c); Section 5.7 introductory text
% Continues from agent_22 which ended mid-sentence at "The remain-"
% =============================================================================

%% ----------------------------------------------------------------
%% 5.6  (continued) -- from page 430
%% ----------------------------------------------------------------

\noindent
ing term, $\mathbf{a} \cdot \mathbf{q}$, is a special relativistic correction associated with the inertia of
the flowing energy $\mathbf{q}$ (see \S\,2.5).

To put the law of energy conservation (5.6.11) into a form relevant to the
radial structure of the disk, (i) we drop the internal energy and compressional
work terms, $\rho_0\, d\Pi/d\tau$ and $\tfrac{1}{3}\theta\, \tilde{p}$, because our condition of ``negligible specific
heat'' (\S\,5.8) guarantees that they are negligible:
%
\begin{align}
&\int \bigl(\rho_0\, d\Pi/d\tau\bigr)\, dz
  \sim \int \bigl(\tfrac{1}{3}\theta\, \tilde{p}\bigr)\, dz \notag\\
&\ll (\text{energy generated by frictional heating}); \notag
\end{align}

(ii) we drop the special relativistic correction $\mathbf{a}\cdot\mathbf{q}$ because the gas is in nearly
geodesic (unaccelerated) orbits; (iii) we average and integrate vertically, and use
the relation $\langle q^{\hat{z}}(r,h)\rangle = \langle -q^{\hat{z}}(r,-h)\rangle = F$;
(iv) we replace the averaged shear of
the gas $\langle\sigma_{\alpha\beta}\rangle$ by the shear of the equatorial geodesic orbits.  The result is
%
\begin{equation}
2F = -\sigma_{\hat{r}\hat{\phi}}^{(\mathrm{EGO})}\,
  \int_{-h}^{h} \langle t^{\hat{r}\hat{\phi}} \rangle\, dz
  = -2\,\sigma_{\hat{r}\hat{\phi}}^{(\mathrm{EGO})}\, W.
\end{equation}
%
Using expression (5.4.6) for $\sigma_{\hat{r}\hat{\phi}}^{(\mathrm{EGO})}$,
we obtain finally
%
\begin{equation}
\tag{5.6.12}
F = \tfrac{1}{2}(Mr^3)^{-1/2}\,\mathscr{G}^{-1/2}\,\mathscr{G}\,W.
\end{equation}

\paragraph*{Manipulation of conservation laws}
The conservation laws (5.6.6) and (5.6.12) for angular momentum and energy,
when combined, give a differential equation for the integrated stress
%
\begin{equation}
\tag{5.6.13}
\bigl(-\dot{M}_0 \tilde{L}/2\pi
  + r^{3/2}\mathscr{G}^{-1/2}\,2\mathscr{G}\,W\bigr)_{,r}
  + \tfrac{1}{2}(Mr)^{-1}\,\tilde{L}\,\mathscr{G}^{-1/2}\,\mathscr{G}\,W = 0.
\end{equation}

All quantities in this differential equation are explicit functions of $r$ that
describe properties of the Kerr metric (\S\,5.4), except the constant accretion
rate $\dot{M}_0$ and the unknown integrated stress $W(r)$.  It is straightforward to integrate
this differential equation.  The constant of integration is fixed by the physical
fact that once the gas reaches the stable circular orbit of minimum radius,
$r = r_{ms}$, the gas will ``fall out'' of the disk and spiral rapidly down the hole.
Consequently, the gas density at $r < r_{ms}$ is virtually zero compared to that at
$r > r_{ms}$---which means that no viscous stresses can act across the surface $r = r_{ms}$;
i.e., $W$ must vanish at $r = r_{ms}$.  The solution that satisfies this boundary
condition is
%
\begin{equation}
\tag{5.6.14a}
W = \frac{\dot{M}_0}{2\pi}\left(\frac{M}{r^3}\right)^{\!1/2}
  \frac{\mathscr{Q}^{1/2}}{\mathscr{B}\,\mathscr{G}}.
\end{equation}

The corresponding value of the flux, as obtained from the law of energy conserva\-tion
(5.6.12), is
%
\begin{equation}
\tag{5.6.14b}
F = \frac{3\dot{M}_0}{8\pi r^2}\,\frac{M}{r}\,
  \frac{\mathscr{Q}}{\tilde{r}\,\mathscr{B}\,\mathscr{G}^{1/2}}.
\end{equation}

The final equations of radial structure are these two equations for $W$ and $F$,
plus the law of rest-mass conservation (5.6.2):
%
\begin{equation}
\tag{5.6.14c}
\dot{M}_0 = -2\pi r\,\Sigma\,\bar{v}^{\hat{r}}\,\mathscr{G}^{1/2}.
\end{equation}

These are the relativistic, black-hole versions of the Newtonian equations of
radial structure (5.2.17), (5.2.18), and (5.2.19).

The equations of radial structure can be viewed as three equations linking
the four radial functions $\Sigma$ (surface density), $W$ (integrated stress), $\bar{v}^{\hat{r}}$ (mass-averaged radial velocity), and $F$ (radiant flux).  Note that these equations
determine $W$ and $F$ explicitly; but they determine only the product $\Sigma\bar{v}^{\hat{r}}$, not
the individual functions $\Sigma$ and $\bar{v}^{\hat{r}}$ individually---and to
calculate other features of the disk such as thickness $2h$, internal temperature
$T$, etc.---one must build a model for the vertical structure.


%% ----------------------------------------------------------------
%% 5.7  Equations of Vertical Structure
%% ----------------------------------------------------------------

\subsection{Equations of Vertical Structure}
\label{sec:5.7}

The equations of radial structure (5.6.14) are based on conservation laws,
without any reference to the detailed properties of the gas in the disk (no
reference to equation of state, or to nature of viscous stresses, or to nature
of opacity, or turbulence and magnetic fields).  By contrast, the
equations of vertical structure, discussed below, require explicit assumptions
about the properties of the gas.  Hence, almost all of the uncertainties and
complications of the model are lumped into the vertical structure.  Ten years
hence one will have a much improved theory of the vertical structure, whereas
the equations of (averaged, steady-state) radial structure will presumably be
unchanged.

Inside the disk the characteristic scale on which the vertical structure changes
is $h$, while the characteristic scale for changes in the radial structure is $r \gg h$.
Consequently, with good accuracy an observer inside the disk can regard the
local variables (temperature, density, etc.)\ as functions of height, $z$, only.  To
analyze the local vertical structure most conveniently, one performs calculations
in the local orbiting orthonormal frame $\mathbf{e}_{\hat{r}}$, $\mathbf{e}_{\hat{\theta}}$,
$\mathbf{e}_{\hat{\phi}}$, which is located in the
central plane of the disk ($z = 0$).  Aside from rotation about the $\mathbf{e}_z$ axis---which
produces Coriolis and centrifugal forces---this frame is (locally) inertial.  The
Coriolis and centrifugal forces have no influence on the vertical structure;
hence we can ignore them and throughout this section can regard the orbiting
frame as inertial.

The laws of physics in the orbiting inertial frame at any given $t$, $r$, $\varphi$ are those
